
%% ═══════════════════════════════════════════════════════════════════════
%%  AUTO-GENERATED by generate_tuning_latex.py
%%  Date: 2026-03-23 20:32
%%  Sources: DQ-MPCC_baseline/tuning/best_weights.json
%%           MPCC_baseline/tuning/best_weights.json
%%
%%  DO NOT EDIT — re-run the script to regenerate.
%% ═══════════════════════════════════════════════════════════════════════

%% ─────────────────────────────────────────────────────────────────────
%%  Bilevel formulation
%% ─────────────────────────────────────────────────────────────────────
\subsubsection{Multi-Velocity Bilevel Problem}
\label{sec:bilevel_problem}

We seek the weight vector
$\boldsymbol{\vartheta}\!\in\!\mathbb{R}^{17}$
that minimises a scalar meta-objective
averaged over a representative velocity set
$\mathcal{V}=\{5, 10, 16\}$~m/s:
%
\begin{equation}
  \boldsymbol{\vartheta}^{\star}
  \;=\;
  \arg\min_{\boldsymbol{\vartheta}\,\in\,\Theta}\;
  J_{\mathrm{multi}}(\boldsymbol{\vartheta}),
  \qquad
  J_{\mathrm{multi}}
  \;=\;
  \frac{1}{|\mathcal{V}|}
  \sum_{v\in\mathcal{V}} J(\boldsymbol{\vartheta},\,v),
  \label{eq:bilevel_multi}
\end{equation}
%
where $\Theta$ is a log-uniform hyper-rectangle (Table~\ref{tab:search_space})
and $J(\boldsymbol{\vartheta},v)$ is the \emph{single-velocity sub-objective}
obtained by running a full closed-loop simulation at
$v_\theta^{\max}=v$.

Each sub-objective aggregates the accumulated MPCC stage cost with
five penalty terms that prevent degenerate solutions
(velocity starvation, incomplete laps, axis-selective tracking):
%
\begin{equation}
\begin{aligned}
J(\boldsymbol{\vartheta},v)
&= \bar{\ell}
+\underbrace{W_{\mathrm{inc}}\,\max(0,\,1\!-\!c)
  \vphantom{\frac{t}{T}}}_{\text{completion}}
+\underbrace{W_{t}\,\frac{t_{\mathrm{lap}}}{T_{\mathrm{ref}}}}_{\text{lap time}}
\\[2pt]
&\quad
+\underbrace{W_{v}\,\frac{[v - \bar{v}_\theta]^{+}}{v}}_{\text{velocity}}
+\underbrace{W_{\alpha}\!\left(\frac{\max_i r_i}{\min_i r_i}-1\right)}_{\text{isotropy}}
\\[2pt]
&\quad
+\underbrace{W_{c}\,\mathrm{RMSE}_{c}\vphantom{\frac{t}{T}}}_{\text{contouring}}\,.
\end{aligned}
\label{eq:sub_objective}
\end{equation}
%
Here
$\bar{\ell}$ is the mean MPCC stage cost over $K$ steps,
$c=\theta_K/s_{\max}$ is the fractional path completion,
$T_{\mathrm{ref}}=s_{\max}/v$ is the ideal lap time,
$[\cdot]^{+}=\max(\cdot,0)$,
$r_i = \bigl(\frac{1}{K}\sum_{k} e_{c,i,k}^{2}+e_{\ell,i,k}^{2}\bigr)^{1/2}$
is the per-axis combined tracking RMSE ($i\!\in\!\{x,y,z\}$),
and $\mathrm{RMSE}_{c}$ is the scalar contouring RMSE.
The fixed penalty weights are
$W_{\mathrm{inc}}\!=\!1000$,
$W_t\!=\!0.5$,
$W_v\!=\!2.0$,
$W_\alpha\!=\!0.3$,
$W_c\!=\!3.0$.

%% ── Search-space table ──
\begin{table}[!t]
\centering
\caption{Log-uniform search space $\Theta$ shared by both controllers
($i\!\in\!\{x,y,z\}$, $j\!\in\!\{x,y,z\}$).
All 17 parameters are sampled independently.}
\label{tab:search_space}
\renewcommand{\arraystretch}{1.1}
\setlength{\tabcolsep}{6pt}
\begin{tabular}{l c}
\toprule
\textbf{Parameter} & \textbf{Range} \\
\midrule
    $Q_{c,i}$            & $[1.0,\;30]$ \\
$Q_{\ell,i}$         & $[0.5,\;30]$ \\
$Q_{{\phi,i}}$       & $[0.1,\;20]$ \\
$U^f$                & $[0.005,\;0.5]$ \\
$U^{{\tau_j}}$       & $[20,\;800]$ \\
$Q_{{\omega,i}}$     & $[0.01,\;2.0]$ \\
$Q_s$                & $[0.5,\;20]$ \\
\bottomrule
\end{tabular}
\end{table}

The outer optimiser is the Tree-structured Parzen Estimator (TPE)
implemented in Optuna~\cite{akiba2019optuna},
with 10 uniform warm-up trials
and seed 42 for reproducibility.
Each trial requires 3 closed-loop simulations
(one per $v\!\in\!\mathcal{V}$), making the per-trial cost
approximately 3$\times$ a single run.

%% ─────────────────────────────────────────────────────────────────────
%%  Tuning results
%% ─────────────────────────────────────────────────────────────────────
\subsubsection{Tuning Results}
\label{sec:tuning_results}

The procedure is executed independently for
both controllers using 100 trials.
The search converges within approximately 72 trials
for both controllers
(total wall time: 19~min for DQ-MPCC,
12~min for the Baseline,
at 11.6~s and 7.3~s per trial respectively).

%% ─────────────────────────────────────────────────────────────────────
%%  TABLE — Optimised cost weights
%% ─────────────────────────────────────────────────────────────────────

\begin{table}[!t]
\centering
\caption{Optimised cost weights obtained by bilevel Bayesian
optimisation (100 TPE trials per controller).
$\bm{q}_\phi$ denotes the $\mathrm{SE}(3)$ rotational penalty (DQ-MPCC)
and $\bm{q}_q$ the $\mathrm{SO}(3)$ rotational penalty (Baseline).}
\label{tab:cost_weights}
\renewcommand{\arraystretch}{1.15}
\setlength{\tabcolsep}{4pt}
\begin{tabular}{l c c}
\toprule
\textbf{Weight} & \textbf{DQ-MPCC} & \textbf{Baseline MPCC} \\
\midrule
    $\bm{q}_{\phi}$ / $\bm{q}_{q}$ & $[2.0, 13.4, 0.66]$ & $[1.8, 0.16, 0.86]$ \\
$\bm{q}_{c}$ & $[12.6, 3.8, 3.3]$ & $[2.5, 15.8, 7.0]$ \\
$\bm{q}_{\ell}$ & $[2.3, 16.8, 1.0]$ & $[26.3, 1.2, 9.6]$ \\
$U^f$ & $8.8\!\times\!10^{-3}$ & $5.7\!\times\!10^{-3}$ \\
$[U^{\tau_x}, U^{\tau_y}, U^{\tau_z}]$ & $[24.1, 72.9, 32.7]$ & $[509, 20.1, 56.1]$ \\
$\bm{q}_{\omega}$ & $[0.040, 0.032, 0.045]$ & $[0.21, 0.020, 0.051]$ \\
$Q_s$ & $0.50$ & $0.58$ \\
\midrule
$J_{\mathrm{multi}}^{\star}$ & $22.18$ & $26.56$ \\
Best trial & \#87 & \#72 \\
\bottomrule
\end{tabular}
\end{table}

Table~\ref{tab:cost_weights} reports the optimal weights.
The DQ-MPCC achieves a $17\%$ lower aggregate objective ($J_{\mathrm{multi}}^\star = 22.18$ vs.\ $26.56$),
indicating that the $\mathrm{SE}(3)$ formulation yields a cost landscape more amenable to Bayesian optimisation.
Both controllers allocate the largest tracking penalties
to the lateral axes ($Q_{c}^x = 12.6$ for DQ-MPCC,
$Q_{c}^y = 15.8$ for the Baseline),
consistent with the Lissajous geometry that generates the
largest contouring excursions in the $y$-axis.
A notable difference is the torque penalty $\bm{w}_u$:
the Baseline requires
$U^{\tau_x} = 509.3$,
roughly 21$\times$ larger
than the DQ-MPCC value
($24.1$),
suggesting that the decoupled Euclidean formulation needs stronger
input regularisation to maintain smooth torque profiles.

